\section{Implications for Design}\label{sec:design}

The framework is meant to be used by anyone who builds autonomous systems. This section draws its implications for one kind of system: a general-purpose coding agent configured to build an application, or to run a study, without a person in the loop. The author is building such a system, an open-source plugin for Claude Code with two commands. One builds an application with a builder, an independent QA reviewer and a tester who uses the running application; the other runs a pre-registered, reproducible study from a question and a dataset. The principles in Table~\ref{tab:design} are stated for any system of this kind, and each carries the confidence of the evidence behind it, not of the project.

\begin{table}[t]
\centering
\footnotesize
\caption{Design principles for unattended systems that build software or run studies, the evidence behind each, and the confidence that evidence supports.}
\label{tab:design}
\begin{tabularx}{\linewidth}{@{}>{\raggedright\arraybackslash}X>{\raggedright\arraybackslash}p{4.2cm}l@{}}
\toprule
Design principle & Evidence & Confidence \\
\midrule
Decide ``done'' with an external signal the builder cannot edit: specification-derived tests and state checks in the running application, run outside the builder's sandbox & \citep{V2,V28,C12}; harbor$^{a}$ & Strong \\
\addlinespace
Tell the builder how it will be checked and let it run its own application, but keep the final check out of its reach & \citep{C29,C9,C19,V3} & Moderate \\
\addlinespace
Treat a completion report as a claim; require the artifacts of execution, such as test output and observed state & \citep{T23,F20,F25} & Strong \\
\addlinespace
If a model reviews the work, have it check observed state against an explicit rubric, not read the builder's narrative & \citep{T18,V24,V29,V15} & Moderate \\
\addlinespace
Add a role only when it is an independent check, and compare the team with a budget-matched single agent & \citep{T18,T25,T9,T4,T6} & Moderate \\
\addlinespace
Isolate parallel builders and merge their work only through executed tests & \citep{T20,T22,T19} & Moderate \\
\addlinespace
Keep the harness lean: few curated skills, no auto-generated context files, deterministic permission rules and a sandbox the agent cannot switch off & \citep{H6,H1,H8,H17}; sandbox-runtime$^{b}$ & Moderate \\
\addlinespace
Bound every run with separate caps and an honest exit, and make a skipped or failed check fail loudly & \citep{V3,S19}; mini-swe-agent$^{c}$; AI-Scientist-v2$^{d}$ & Suggestive \\
\addlinespace
For studies, commit the hypothesis, decision rule and held-out split before results exist; keep the grader and answers out of the agent's workspace; re-run every reported number from code & \citep{R19,R6,R26,S15}; RD-Agent$^{e}$ & Moderate \\
\addlinespace
Report $k$ runs per arm with the spread, cost per passed task and interventions per task, graded by an evaluator that did not build the system & \citep{F19,F24,F8,F32} & Strong \\
\bottomrule
\multicolumn{3}{@{}p{\textwidth}@{}}{\scriptsize\raggedright Repositories: $^{a}$\,\repofile{harbor-framework/harbor}{src/harbor/cli/quality_checker/default-rubric.toml}; $^{b}$\,\repo{anthropics/sandbox-runtime}{97}; $^{c}$\,\repofile{SWE-agent/mini-swe-agent}{src/minisweagent/agents/default.py}; $^{d}$\,\repo{SakanaAI/AI-Scientist-v2}{133}; $^{e}$\,\repofile{microsoft/RD-Agent}{rdagent/scenarios/data_science/test_eval.py}.} \\
\end{tabularx}
\end{table}

The table sorts into three groups. The strongest principles concern who decides that work is done. Checks that execute outside the builder's reach are the one cell of Figure~\ref{fig:twobytwo} that gave reliable evidence, while the instrument itself was sound \citep{V2,V28,C12}, and the evidence that completion reports mislead and that single runs are noisy is among the firmest in the review \citep{F20,T23,F19,F24}. The moderate principles concern how a system is organized: which roles to add, how to run parallel builders, how much harness to carry and how to keep an agent-run study honest. There the direction of the evidence is consistent, but the studies are few, small or written by the teams that built the method. The weakest concerns bounds. Compounding is well documented \citep{S19}, but which caps and exits help is known mainly from one study of honest exits \citep{V3} and from practitioner tools.

Two constraints for a system of this kind come from the platform rather than the literature. A run that must go for hours without a person cannot wait for permission prompts to be approved, so permissions have to come from rules set before launch; the author's design scopes them to the project folder and an allow-list of commands, the kind of deterministic control that Section~\ref{sec:harness} found more reliable than model judgment. And on the author's platform, separate roles run as separate sessions that hand off through files, so the information lost at each handoff (Section~\ref{sec:architecture}) is a risk to measure rather than assume away.

Practitioner playbooks agree with the evidence on some points and not others. Some encode it directly: one skills plugin forbids any completion claim made without fresh verification evidence (\repofile{obra/superpowers}{skills/verification-before-completion/SKILL.md}), which is the lesson of Section~\ref{sec:verification} written as a rule. The team methods reviewed in Section~\ref{sec:architecture} organize work around roles and handoffs and, in our reading, report no measured success rates; MetaGPT's current default team hires no QA engineer and defaults to not running tests (\repofile{FoundationAgents/MetaGPT}{metagpt/software_company.py}), which omits the role the evidence most supports. The popular research loops get the signal right, keeping a change only if a validation number improves (\repofile{karpathy/autoresearch}{program.md}), but leave the call to the evaluator in a file the agent owns (\repo{karpathy/autoresearch}{599}) and select repeatedly on the same validation number, a practice that overfit the validation score in a study of search loops \citep{R26}. The author's own playbook assigns a lead, builders, QA, a tester, a marketing role, an auditor and an evaluator. Of these, QA, the tester and the evaluator are independent checks on the builder's work. The evidence reviewed here gives no accuracy-based reason for the rest, which must be justified, if at all, on other grounds such as coordination. The general rule is that a role-heavy team is justified only where a role is an independent check, and a system of this kind should hold its own claims to the reporting standard of Table~\ref{tab:reporting}.
